\documentclass[10pt,a4paper]{amsart}
\usepackage[margin=19mm]{geometry}
\usepackage{amsmath,amssymb,mathtools}
\usepackage[scr=rsfs,cal=euler]{mathalfa}
\usepackage{xfrac}
\usepackage[unicode,hidelinks,bookmarks=false]{hyperref}
\newcommand{\ie}{i.e.\ }
\newcommand{\cf}{cf.\ }
\newcommand{\resp}{respectively\ }
\newcommand{\Subsection}{\subsection}
\providecommand{\nobreakdash}{\nobreak\hbox{-}}
\setlength{\emergencystretch}{2em}
\multlinegap0pt
\usepackage{amssymb}
\usepackage{mathtools}
\usepackage{tikz}
\usepackage[shortlabels]{enumitem}
\usepackage{xcolor}
\newtheorem{res}{Result}
\newtheorem{lm}{Lemma}
\newtheorem{prop}{Proposition}
\usepackage{cleveref}
\crefname{res}{Result}{Results}
\crefname{lm}{Lemma}{Lemmas}
\crefname{prop}{Proposition}{Propositions}
 \newcommand{\FF}{\mathbb F}
 \newcommand{\eps}{\varepsilon}
 \newcommand{\mc}{\mathcal C}
 \newcommand{\mg}{\mathcal G}
 \newcommand{\mq}{%\mathcal
 Q}
 \newcommand{\mr}{\mathcal R}
\newcommand{\old}[1]{}%\textcolor{blue}{#1}}
 \newcommand{\pg}{\textnormal{PG}}
 \DeclareMathOperator{\proj}{proj}
 \newcommand{\set}[1]{ \left \{ #1 \right \} }
 \newcommand{\sett}[2]{ \left\{ #1 \, \, || \, \, #2 \right \} }
 \newcommand{\vspan}[1]{\left \langle #1 \right \rangle}
 \newcommand{\zero}{\mathbf 0}
\title{The minimum weight of the code of intersecting lines in $\pg(3,q)$}
\author{Sam Adriaensen, Robin Simoens, Leo Storme}
\date{Source: arXiv:2403.07451v1 (CC BY 4.0)}
\begin{document}
\maketitle

\noindent\emph{Source and licence.} The minimum weight of the code of intersecting lines in $\pg(3,q)$. Version arXiv:2403.07451v1 (CC BY 4.0), distributed by arXiv under the Creative Commons Attribution 4.0 International licence, http://creativecommons.org/licenses/by/4.0/, as recorded in the arXiv licence metadata for this exact version and shown on the abstract page https://arxiv.org/abs/2403.07451v1. Full licence notice: \texttt{licenses/arxiv-2403.07451-CC-BY-4.0.txt}.\par\smallskip

\noindent\textit{Editorial scope.} Counted unit: the article's prop Prop:Symplectic (source0.tex lines 672--675). The article's complete original proof of it is delivered with it (source0.tex lines 677--791, one \texttt{proof} environment of 115 lines). No mathematics is written, repaired or re-derived: the statement and the proof are the article's own lines, in the article's own notation, and no retained line is substituted or re-flowed.  Retained uncounted as the source-local context the proof needs: lines 326--350; lines 622--624.  Nothing was omitted from any retained span, so the package carries no omission record.  No retained line contains a citation, so the wrapper carries no bibliography and no bibliography entry.  No retained line contains a figure environment, an image-inclusion command or drawing code, so no image is delivered and the package declares no asset dependency.  Editorial formatting only, in addition: an AMS \texttt{amsart} preamble that restates the article's own declarations and macros used by the retained lines, namely the environments \texttt{res}, \texttt{lm}, \texttt{prop} on separate counters, together with the standard packages those lines need.  The retained environments print in this document as Res 1, Lm 1, Proposition 1 in order of appearance, whereas the article prints the same items with the numbering of its own full document.  The complete native source of the article as distributed in the arXiv e-print for this exact version is bundled unchanged as \texttt{sources/155955/source0.tex}.
Let $\pi$ and $\rho$ be two disjoint subspaces of $\pg(n,q)$ and let $S$ be a set of points in $\pi$.
The set $\rho S = \bigcup_{P \in \pi} \vspan{\rho,P}$ is called the \emph{cone} with \emph{base} $S$ and \emph{vertex} $\rho$.
By convention, $\rho S = S$ if $\rho = \emptyset$ and $\rho S = \rho$ if $S = \emptyset$.
Every quadric is a cone with as base a non-degenerate quadric.
We call a subspace $\pi$ \emph{singular} with respect to a quadric $Q$ if $\pi \cap Q$ is a degenerate quadric.
If $P \in Q^\eps(n,q)$ and $Q^\eps(n,q)$ has a polarity $\perp$, then $P^\perp$ intersects $Q^\eps(n,q)$ in a quadric $P Q^\eps(n-2,q)$.
We call $P^\perp$ the \emph{tangent hyperplane} of $P$.
A subspace contained in $Q^\eps(n,q)$ of maximal dimension is called a \emph{generator}.
\begin{res}
 \label{Res:Quadrics}
 \begin{enumerate}
  \item The quadric $Q^\eps(n,q)$ has $\frac{q^n-1}{q-1} + \eps q^\frac{n-1}2$ points.
  \item \label{Res:Quadrics:Generators} The quadric $Q^\eps(n,q)$ has
  \[
   \prod_{i=1-\eps}^\frac{n-\eps}2 (q^i+1)
  \]
  generators, and their dimension equals $\frac{n+\eps}2-1$.
  \item \label{Res:Quadrics:Equivalence} Let $\mg$ denote the set of generators of $Q^+(n,q)$.
  Then the relation
  \[
   \sett{ (\pi, \rho) \in \mg^2 }{ \dim (\pi \cap \rho) \equiv \frac{n-1}2 \pmod 2 }
  \]
  is an equivalence relation on $\mg$, having two classes of equal size.
 \end{enumerate}
\end{res}

\begin{lm}\label{Lm:Proj}
Let $0<i,j,k<n$, $i<j$. A function $c: \mg_j(n,q) \to \FF_p$ belongs to  $\mc_{j,k}(n,q)^\perp$ if and only if $\proj^{(i)}(c)$ belongs to $\mc_{i,k}(n,q)^\perp$.
\end{lm}

\begin{prop}
 \label{Prop:Symplectic}
 Let $S$ denote the set of absolute lines of a symplectic polar space $W(2n+1,q)$ embedded in $\pg(2n+1,q)$. Then $\chi_S \in \mc_{1,n}(2n+1,q)$ if and only if $q$ is even.
\end{prop}

\begin{proof}
 \old{First suppose that $q$ is even.
 Consider a hyperbolic quadric $\mq^+(3,q)$, and a regulus $\mathcal R$ of lines contained in $\mq^+(3,q)$.
 Consider the codeword 
 \[
  c = \sum_{\ell \in \mathcal R} \chi_\ell
 \]
 of the code $C(q)$.
 Each point of $\mq^+(3,q)$ is contained in a unique line of $\mathcal R$.
 Take a line $\ell$ in $\pg(3,q)$.
 If $\ell$ intersects $\mq^+(3,q)$ in $k \in \set{0,1,2}$ points, then clearly $c\cdot \chi_\ell = k$.
 Otherwise, $\ell$ is completely contained in $\mq^+(3,q)$.
 Either $\ell$ belongs to $\mathcal R$, and $c\cdot \chi_\ell = 1$, or $\ell$ belongs to the opposite regulus and $c\cdot \chi_\ell = q+1 = 1$.
 Hence, $c$ is the characteristic vector of the singular lines of $\mq^+(3,q)$.
 Since $q$ is even, these are the absolute lines with respect to the polarity corresponding to $\mq^+(3,q)$, which must be symplectic, again by $q$ being even.}

 First suppose that $q$ is even.
 Consider a hyperbolic quadric $\mq^+(2n+1,q)$, and an equivalence class $\mathcal R$ of the generators (\(n\)-spaces) of $\mq^+(2n+1,q)$.
 Consider the codeword 
 \[
  c = \sum_{\lambda\in \mathcal R} \chi^{(1)}_{\lambda}
 \]
 of the code $\mc_{1,n}(2n+1,q)$.
 Take a line $\ell$ in $\pg(2n+1,q)$.
 Then $c(\ell)$ equals the number of generators in $\mr$ that intersect $\ell$ non-trivially.
 Suppose that $\ell$ intersects $Q^+(2n+1,q)$ in $m_1$ points.
 The number of generators of $Q^+(2n+1,q)$ of $\mr$ through a point $P$ equals 0 if $P \notin Q^+(2n+1,q)$ and equals the number of generators of $\mr$ in $P^\perp$ otherwise, with $\perp$ the polarity associated to $Q^+(2n+1,q)$.
 Since $P^\perp$ intersects $Q^+(2n+1,q)$ in a cone $P Q^+(2n-1,q)$, this number equals the number of generators of one equivalence class in $Q^+(2n-1,q)$, which is given by $\prod_{i=1}^{n-1} (q^i+1)$ (where we use the convention that an empty product equals $1$) by \Cref{Res:Quadrics} (\ref{Res:Quadrics:Generators}) and (\ref{Res:Quadrics:Equivalence}).
 Let $m_2$ denote the number of generators of $Q^+(2n+1,q)$ through $\ell$.
 Then
 \[
  c(\ell) = m_1 \prod_{i=1}^{n-1} (q^i + 1) - q m_2 \equiv m_1 \pmod p.
 \]
 Hence, $c$ is the characteristic vector of the singular lines of $Q^+(2n+1,q)$.
 Since $q$ is even, these are the absolute lines with respect to $\perp$, which must be symplectic, again by $q$ being even.
 
 \old{Now suppose that $q$ is odd.
 Consider again a hyperbolic quadric $\mq^+(3,q)$, and its two reguli $\mathcal R^+$ and $\mathcal R^-$.
 Consider the function
 \[
  c: \mg_1(3,q) \to \FF_p: \ell \mapsto \begin{cases}
   1 & \text{if } \ell \in \mathcal R^+, \\
   -1 & \text{if } \ell \in \mathcal R^-, \\
   0 & \text{otherwise.}
  \end{cases}
 \]
 Then $\proj^{(0)}(c) = \zero$, hence $c \in \mc(3,q)^\perp$ by \Cref{Lm:Proj}.
 Therefore, if $\chi_S$ would be in $\mc(3,q)$, then $\chi_S \cdot c = 0$, which means that $S$ contains the same number of lines of $\mathcal R^+$ and $\mathcal R^-$ modulo $p$.
 Hence, in order to prove that $\chi_S\not\in C(q)$, it suffices to show that there exists a hyperbolic quadric $Q^+(3,q)$ and a symplectic polar space $W(3,q)$ in $\pg(3,q)$ where the reguli of $Q^+(3,q)$ do not contain the same number of lines of $W(3,q)$ modulo $p$.
 To this end, consider the symplectic polar space defined by the bilinear form
 \[
  b(X,Y) = X_0 Y_1 - X_1 Y_0 + X_2 Y_3 - X_3 Y_2.
 \]
 A line through two distinct points of $\pg(3,q)$, with coordinates $X=(X_0,X_1,X_2,X_3)$ and $Y=(Y_0,Y_1,Y_2,Y_3)$, is absolute if and only if $b(X,Y) = 0$.
 Now consider the hyperbolic quadric with equation $X_0 X_1 = X_2 X_3$.
 Its two reguli are
 \begin{align*}
  \mathcal R^+ &= \sett{\vspan{(\alpha,0,\beta,0), (0,\beta,0,\alpha)}}{ \vspan{(\alpha,\beta)} \in \pg(1,q)}, \\
  \mathcal R^- &= \sett{\vspan{(\alpha,0,0,\beta), (0,\beta,\alpha,0)}}{\vspan{(\alpha,\beta)} \in \pg(1,q)}.
 \end{align*}
 On the one hand,
 \[
  b\left((\alpha,0,\beta,0),(0,\beta,0,\alpha)\right)
  = 2 \alpha \beta.
 \]
 Since $q$ is odd, this means that $\mathcal R^+$ contains two absolute lines.
 On the other hand
 \[
  b\left((\alpha,0,0,\beta),(0,\beta,\alpha,0)\right) = \alpha \beta - \beta \alpha=0,
 \]
 so all $q+1$ lines in $\mathcal R^-$ are absolute.
 Note that $q+1 \equiv 1 \not \equiv 2 \pmod p$, which finishes the proof.}

 Now suppose that $q$ is odd.
 Let $\perp$ denote the polarity associated to $W(2n+1,q)$.
 Take a 3-space $\Sigma$ of $W(2n+1,q)$ such that $\Sigma \cap \Sigma^\perp = \emptyset$.
 Then the absolute subspaces of $\perp$ contained in $\Sigma$ are the lines and points of a $W(3,q)$ in $\Sigma$.
 Consider a hyperbolic quadric $Q^+(3,q)$ in $\Sigma$ and let $\mr^+$ and $\mr^-$ denote its two reguli.
 Consider the function
 \[
  c: \mg_1(2n+1,q) \to \FF_p: \ell \mapsto \begin{cases}
   1 & \text{if } \ell \in \mathcal R^+, \\
   -1 & \text{if } \ell \in \mathcal R^-, \\
   0 & \text{otherwise.}
  \end{cases}
 \]
 Then $\proj^{(0)}(c) = \zero$, hence $c \in \mc_{1,n}(2n+1,q)^\perp$ by \Cref{Lm:Proj}.
 Therefore, if $\chi_S$ would be in $\mc_{1,n}(2n+1,q)$, then $\chi_S \cdot c = 0$, which means that $S$ contains the same number of lines of $\mathcal R^+$ and $\mathcal R^-$ modulo $p$.
 Hence, in order to prove that $\chi_S \notin \mc_{1,n}(2n+1,q)$, it suffices to show that there exists a hyperbolic quadric $Q^+(3,q)$ and a symplectic polar space $W(3,q)$ in $\Sigma \cong \pg(3,q)$ where the reguli of $Q^+(3,q)$ do not contain the same number of lines of $W(3,q)$ modulo $p$.
 
 To this end, consider the symplectic polar space of $\pg(3,q)$ defined by the bilinear form
 \[
  b(x,y) = x_0 y_1 - x_1 y_0 + x_2 y_3 - x_3 y_2.
 \]
 From now on, let $\perp$ denote the polarity of the symplectic space $W(3,q)$ corresponding to $b$.
 A line through two distinct points of $\pg(3,q)$, with coordinates $x=(x_0,\dots,x_3)$ and $y=(y_0,\dots,y_3)$, is absolute if and only if $b(x,y) = 0$.
 Now consider the hyperbolic quadric with equation $x_0 x_1 = x_2 x_3$.
 Its two reguli are
 \begin{align*}
  \mathcal R^+ &= \sett{\vspan{(\alpha,0,\beta,0), (0,\beta,0,\alpha)}}{ \vspan{(\alpha,\beta)} \in \pg(1,q)}, \\
  \mathcal R^- &= \sett{\vspan{(\alpha,0,0,\beta), (0,\beta,\alpha,0)}}{\vspan{(\alpha,\beta)} \in \pg(1,q)}.
 \end{align*}
 On the one hand,
 \[
  b\left((\alpha,0,\beta,0),(0,\beta,0,\alpha)\right)
  = 2 \alpha \beta.
 \]
 Since $q$ is odd, this means that $\mathcal R^+$ contains two absolute lines.
 On the other hand
 \[
  b\left((\alpha,0,0,\beta),(0,\beta,\alpha,0)\right) = \alpha \beta - \beta \alpha=0,
 \]
 so all $q+1$ lines in $\mathcal R^-$ are absolute.
 Note that $q+1 \equiv 1 \not \equiv 2 \pmod p$, which finishes the proof.
\end{proof}

\end{document}
